\chapter{Alphabete, Wörter, Sprachen}
Wenn man sich mit der Funktionsweise von Rechnern (Computern) genauer beschäftigt, stellt man fest, dass Rechner im Grunde eine Transformation von Eingabedaten in Ausgabedaten realisieren. Sowohl die Eingabedaten als auch die Ausgabedaten lassen sich als Texte darstellen. Die Texte sind nichts anderes als Folgen von Symbolen aus einem bestimmten Alphabet. Programme können als Folge von Symbolen der Computertastatur dargestellt werden. In digitalen Rechnern sind alle Informationen als Folgen von Einsen und Nullen gespeichert. Damit realisiert der Rechner eine Transformation von Eingabetexten in Ausgabetexte.\\

\noindent
In diesem Kapitel wollen wir den Formalismus für den Umgang mit Texten kennenlernen. Wir werden die fundamentalen Begriffe \textit{Alphabet}, \textit{Wort} und S\textit{prache} einführen. Diese werden uns später helfen, bekannte Probleme der Informatik wie beispielsweise das \textit{Entscheidungsproblem} mathematisch sauber zu formulieren.\\

\noindent
Dieses Kapitel folgt dem Abschnitt 2.2 aus dem Buch\footnote{J. Hromkovic: Theoretische Informatik. 5. Auflage, Springer Vieweg 2014., ISBN: 978-3-658-06432-7} sehr  nahe. Es wurden einige Aufgaben, welche als Hilfestellung dienen, hinzugefügt und kleine Teile ausgelassen.

\section{Alphabete, Wörter, Sprachen}
\begin{deff}[Alphabet]{deff:Alphabet}
Eine \textbf{endliche nichtleere} Menge $\Sigma$ heisst \textit{Alphabet}. Die \textit{Elemente eines Alphabets} werden \textit{Buchstaben} (\textit{Zeichen, Symbole}) genannt.
\end{deff}
Wir werden später Alphabete verwenden, um eine schriftliche Darstellung einer Sprache zu erzeugen.\\
Die Definition, gegeben in \ref{deff:Alphabet}, entspricht unserer intuitiven Vorstellung eines Alphabets:\\
Um Text darstellen zu können, muss ein Alphabet mindestens ein Symbol enthalten ($\to$ nichtleer). Damit man sich auf einen Satz von Zeichen einigen kann, darf das Alphabet nicht unendlich gross sein ($\to$ endlich).\\
\clearpage
\noindent
Wir listen nun einige der Alphabete auf, die in der Mathematik und Informatik häufig verwendet werden:
\begin{itemize}
    \item $\Sigma_{\text{bool}} = \menge{0,1}$ ist das Boole'sche Alphabet, mit dem digitale Rechner arbeiten.
    \item $\Sigma_{\text{lat}} = \menge{a,b,c,...,z,A,B,...,Z}$ ist das lateinische Alphabet.
    \item $\Sigma_{\text{Tastatur}} = \menge{a,b,c,...,z,A,B,C,...,Z,\text{\textvisiblespace}\,,>,<,(,),...,\#,?,!}$ ist das Alphabet aller Symbole, die mit der englischen Tastatur getippt werden können. Dabei ist \textvisiblespace\ das Symbol für das Leerzeichen / Leersymbol.
    \item $\Sigma_{\text{greek}} = \menge{\alpha, \beta, \gamma, ..., \omega, A, B, \Gamma, ..., \Omega}$ ist das griechische Alphabet.
    \item $\Sigma_m = \menge{0,1,2,...,m-1}$ für jedes $m\in\mathbb{N}-\menge{0}$ ist ein Alphabet für die $m$-adische Darstellung von Zahlen.
\end{itemize}
\textit{Wörter} werden wir als endliche Folgen von Buchstaben ansehen.
\begin{deff}[Wort]{deff:Wort}
Sei $\Sigma$ ein Alphabet. Ein \textit{Wort} über $\Sigma$ ist eine endliche (möglicherweise leere) Folge von Buchstaben aus $\Sigma$. Das \textit{leere Wort} $\lambda$ ist die leere Buchstabenfolge.
\end{deff}
Die \textit{Länge} $|w|$ eines Wortes $w$ ist die Länge des Wortes als Folge, das heisst die Anzahl der Vorkommen von Buchstaben in $w$.\\

\noindent
\underline{Beispiele}:\\
$w = 1,0,0,1,0$ ist ein Wort über dem Alphabet $\Sigma_{\text{bool}}$ und $|w| = 5$, da $w$ eine Folge von $5$ Buchstaben ist.\\
$u = M,\text{\textvisiblespace}\,, j$ ist ein Wort über $\Sigma_{\text{Tastatur}}$ und $|u| = 3$, da $u$ eine Folge von $3$ Buchstaben ist.\\
Das leere Wort $\lambda$ ist ein Wort über jedem Alphabet und es gilt $|\lambda| = 0$.
\begin{Exercise}%[difficulty={1}]
\noindent
Was ist $\Sigma_{10}$?\\
\vspace{0.6cm}
\end{Exercise}
\begin{Answer}
\noindent
\begin{align*}
    \Sigma_{10} = \menge{0,1,2,...,9}
\end{align*}
ist das Alphabet für die Darstellung von Zahlen im Dezimalsystem.
\end{Answer}

\noindent
Sei $\Sigma$ ein Alphabet. $\Sigma^*$ (gesprochen: \textit{Sigma Stern}) ist \textbf{die Menge aller Wörter} über $\Sigma$, also die Menge aller endlichen Folgen von Symbolen aus $\Sigma$.\\
$\Sigma^+ = \Sigma^*-\menge{\lambda}$ ist die Menge aller Wörter über $\Sigma$ ohne das leere Wort.\\

\noindent
Wir werden im Folgenden Wörter ohne Kommma schreiben. Anstelle von $1,0,0,1,0$ werden wir lediglich $10010$ schreiben. Allgemein, werden wir anstelle von $x_1,x_2,...,x_n$ einfach $x_1x_2...x_n$ schreiben.
\clearpage

\noindent
\underline{Bemerkung}:\\
In der deutschen Sprache existieren die zwei Ausdrücke \quotes{Wörter} und \quotes{Worte}. Dies sind keine Synonyme. \textit{Wörter} bezeichnet den Plural von \textit{Wort} (\quotes{Im Duden stehen viele Wörter}). Der Ausdruck \textit{Worte} bezieht sich auf Gedankenkonstrukte (\quotes{Sie sprach weise Worte}).\\

\noindent
Wir können Wörter benutzen, um mathematische Objekte wie Zahlen, Formeln, Graphen und Computerprogramme darzustellen.\\
Ein Wort $x = x_1x_2...x_n\in (\Sigma_{\text{bool}})^*$ kann als binäre Darstellung der Zahl
\begin{align}\label{eq:Nummer}
    \text{Nummer}(x)=\displaystyle\sum_{i=1}^n x_i\cdot 2^{n-i}
\end{align}
betrachtet werden.

\begin{Exercise}%[difficulty={1}]
\noindent
Sei $x = 1011$. Berechnen Sie $\text{Nummer}(x)$.\\
\vspace{0.6cm}
\end{Exercise}
\begin{Answer}
\noindent
\begin{align*}
    \text{Nummer}(x) = \text{Nummer}(\red{1}\blue{0}\green{1}\gray{1}) = \displaystyle\sum_{i=1}^4 x_i\cdot 2^{4-i} = \red{1}\cdot 2^3+\blue{0}\cdot 2^2 + \green{1}\cdot 2^1+\gray{1}\cdot 2^0 = 8 + 2 + 1 = 11
\end{align*}
\end{Answer}
\noindent
Für eine Zahl $m\in\mathbb{N}-\menge{0}$ wird mit $\text{Bin}(m)\in (\Sigma_{\text{bool}})^*$ die kürzeste binäre Darstellung von $m$ bezeichnet, also gilt $\text{Nummer}(\text{Bin}(m)) = m$. Man setzt $\text{Bin}(0) = 0$.

\begin{Exercise}%[difficulty={1}]
\noindent
Sei $m\in\mathbb{N}-\menge{0}$. Welche Bedingung muss gelten, damit $\text{Bin}(m)$ die \textbf{kürzeste} binäre Darstellung von $m$ ist.\\
\end{Exercise}
\begin{Answer}
\noindent
Die Bedingung ist lediglich, dass die Darstellung mit dem Symbol $1$ beginnt, da führende Nullen die Bedeutung der Darstellung nicht verändern, sondern die Darstellung lediglich länger machen. Zum Beispiel beschreiben $101$, $0101$ und $000000101$ dieselbe Zahl. Die erste ist aber die kürzeste Darstellung.
\end{Answer}

\begin{theo}[Länge der binären Darstellung]{theo:laenge_binaer}
Sei $m\in\mathbb{N}-\menge{0}$. Dann ist die Länge $|\text{Bin}(m)|$ der kürzesten binären Darstellung von $m$ gegeben durch
\begin{align*}
    |\text{Bin}(m)| = \ceil{\log_2(m+1)}
\end{align*}
\end{theo}
Dabei bezeichnet für eine reelle Zahl $x$ der Ausdruck $\ceil{x}$ die nächste ganze Zahl $\geq x$. Genauer:\\
Sei $x$ eine reelles Zahl, dann definiert man
\begin{align*}
    \ceil{x} = \min\mengecm{k\in\mathbb{Z}}{k\geq x}
\end{align*}
und analog
\begin{align*}
    \floor{x} = \max\mengecm{k\in\mathbb{Z}}{k\leq x}.
\end{align*}
\begin{Exercise}%[difficulty={1}]
\noindent
Bestimmen Sie mit Hilfe der Formel aus Satz \ref{theo:laenge_binaer} die Länge der kürzesten binären Darstellung der Zahl 58902.
\end{Exercise}
\begin{Answer}
\noindent
Mit Hilfe eines Taschenrechners finden wir
\begin{align*}
    \log_2(58902 + 1)\approx 15.846
\end{align*}
und damit
\begin{align*}
    \ceil{\log_2(58902 + 1)} = 16
\end{align*}

\end{Answer}

\clearpage
\begin{Exercise}%[difficulty={1}]
\noindent
Beweisen Sie Satz \ref{theo:laenge_binaer}.\\
Tipp:
\begin{itemize}
    \item Setzen Sie $k = |\text{Bin}(m)|$.
    \item Erklären Sie, warum die Sandwich Beziehung $2^k-1 \geq m > 2^{k-1} - 1$ gilt, indem Sie folgende Fragen beantworten:
    \begin{itemize}
        \item Was ist der Dezimalwert der kleinsten binären Zahl der Länge $k+1$ (ohne führende Nullen)?
        \item Was ist der Dezimalwert der kleinsten binären Zahl der Länge $k$ (ohne führende Nullen)?
    \end{itemize}
\end{itemize}
\end{Exercise}
\begin{Answer}
\noindent
Sei $k = |\text{Bin}(m)|$. Wir wollen $k$ ermitteln.\\

\noindent
Die kleinste binäre Zahl der Länge $k+1$ ist $1\underbrace{00...0}_{k\;\text{mal}}$. Diese hat gemäss der Gleichung \ref{eq:Nummer} den dezimalen Wert (Nummer) $2^{k}$. Dann kann der Dezimalwert der grössten binäre Zahl der Länge $k$ höchstens $2^{k-1}-1$ sein ($1$ weniger). Damit ist $2^k-1 \geq m$ begründet.\\

\noindent
Die kleinste binäre Zahl der Länge $k$ ist $1\underbrace{00...0}_{(k-1)\;\text{mal}}$, was dem dezimalen Wert $2^{k-1}$ entspricht. Also muss auch $m > 2^{k-1} - 1$ gelten, was die Sandwich-Beziehung aus dem Tipp begründet.\\

\noindent
Offensichtlich gilt nun also für $m$ die Beziehung
\begin{align*}
    2^k-1 \geq m > 2^{k-1} - 1 \Rightarrow \\
    2^k \geq m+1 > 2^{k-1} \Rightarrow \\
    k \geq \log_2(m+1) > k-1.
\end{align*}
Die Beziehung $k \geq \log_2(m+1) > k-1$ sagt aus, dass $\log_2(m+1)\in (k-1,k]$ und damit folgt:
\begin{align*}
    k = \ceil{\log_2(m+1)}.
\end{align*}
\end{Answer}

\begin{Exercise}\label{exercise:2.6}%[difficulty={1}]
\noindent
Wie sieht die Menge $\menge{1}^*$ aus? Und wie die Menge $(\Sigma_{\text{bool}})^*$?\\
\end{Exercise}
\begin{Answer}
\noindent
\begin{align*}
    \menge{1}^* &= \menge{\lambda, 1, 11, 111, 1111, 11111, ...} =\\
    &= \menge{\lambda}\cup \mengecm{x_1x_2...x_i}{i\in\mathbb{N},\; x_j = 1\;\text{für}\; j = 1,2,..., i}
\end{align*}
\begin{align*}
    (\Sigma_{\text{bool}})^* = \menge{0,1}^* &= \menge{\lambda, 0, 1,, 00, 01, 10, 11, 000, 001, 010, 100, 011, ...} =\\
    &= \menge{\lambda}\cup \mengecm{x_1x_2...x_i}{i\in\mathbb{N},\; x_j\in\Sigma_{\text{bool}} \;\text{für}\; j = 1,2,..., i}
\end{align*}
\end{Answer}
\vspace{1cm}

\begin{Exercise}%[difficulty={1}]
\noindent
In Definition \ref{deff:Wort}, haben wir ein Wort als endliche Folge von Buchstaben definiert. Nun enthält zum Beispiel die Menge $\menge{1}^*$ in Übung \ref{exercise:2.6} aber Wörter beliebiger Länge. Erklären Sie, warum dies kein Widerspruch ist.
\end{Exercise}
\begin{Answer}
\noindent
Anhand dieser Frage lässt sich der Unterschied zwischen den Begriffen \textit{unbeschränkt} und \textit{unendlich} schön verdeutlichen. Betrachten wir dazu exemplarisch nochmals die Menge $menge{1}^*$ aus Übung \ref{exercise:2.6}:
\begin{align*}
    \menge{1}^* &= \menge{\lambda, 1, 11, 111, 1111, 11111, ...} =\\
    &= \menge{\lambda}\cup \mengecm{x_1x_2...x_i}{i\in\mathbb{N},\; x_j = 1\;\text{für}\; j = 1,2,..., i}
\end{align*}
\begin{itemize}
    \item Für jede natürliche Zahl $n\in\mathbb{N}$, existiert in $\menge{1}^*$ ein Wort $w$ mit $|w| \geq n$, da für jedes $n\in\mathbb{N}$ (insbesondere) auch das Wort $w:=1^n$ in der Menge $\menge{1}^*$ enthalten ist und $|w|\geq n$. Damit gibt es \textbf{keine obere Schranke} für die Länge der Wörter in $\menge{1}^*$.
    \item Jede Wort $w\in\menge{1}^*$ hat die Form $w = 1^n$ für eine natürliche Zahl $n$. Damit hat jedes Wort in der Menge $\menge{1}^*$ eine endliche Länge.
\end{itemize}
Die Länge der Wörter in $\menge{1}^*$ können also nicht nach oben beschränkt werden, trotzdem hat jedes Wort in $\menge{1}^*$ eine endliche Länge und ist somit tatsächlich ein Wort im Sinne von Definition \ref{deff:Wort}.
\end{Answer}
\vspace{1cm}
\noindent
Da Alphabete insbesondere auch Mengen sind, kann man auch das kartesische Produkt von Alphabeten bilden.

\noindent
\underline{Beispiel}:\\
Sei $\Sigma_1 = \menge{g,h}$ und $\Sigma_1 = \menge{x,y}$, dann ist das kartesische Produkt $\Sigma_1\times \Sigma_2 = \menge{(g,x),(g,y),(h,x),(h,y)}$.\\

\noindent
Wir werden nun eine Operation einführen, welche uns erlaubt Wörter zu verketten. Diese Operation werden wir sehr häufig verwenden.
\begin{deff}[Verkettung / Konkatenation]{deff:Konkat}
Die \textbf{Verkettung} (\textbf{Konkatenation)}) für ein Alphabet $\Sigma$ ist eine Abbildung\\
Kon: $\Sigma^*\times \Sigma^*\rightarrow\Sigma^*$, so dass
\begin{align*}
    \text{Kon}(x,y) = x\cdot y = xy
\end{align*}
für alle $x,y\in\Sigma^*$.
\end{deff}
\noindent
\underline{Beispiel}:\\
Sei $\Sigma = {8,9,d,e}$ und seien $x = e899dd$ und $y = de8$, dann ist $\text{Kon}(x,y) = x\cdot y = e899ddde8$.\\

\noindent
Wir werden fast ausnahmslos $xy$ schreiben und nur selten $x\cdot y$ oder $\text{Kon}(x,y)$.\\

\begin{Exercise}%[difficulty={1}]
\noindent
Erläutern Sie in Ihren eigenen Worten, was die Bedeutung des Ausdrucks $\Sigma^*\times \Sigma^*\rightarrow\Sigma^*$ in der Definition \ref{deff:Konkat} der Verkettung ist.\\
\end{Exercise}
\begin{Answer}
\noindent
Wir haben zwei Wörter $x$ und $y$ über dem Alphabet $\Sigma$. Weil $\Sigma^*$ die Menge aller Wörter über $\Sigma$ ist, gilt offensichtlich $x,y\in\Sigma^*$.\\
Aus den zwei Wörtern $x$ und $y$ wird nun ein neues Wort $xy$ gebildet. Man beachte, dass die Reihenfolge $xy\neq yx$ (im Allgemeinen) eine Rolle spielt.\\
Die Verkettung nimmt also ein geordnetes Paar $(x,y)\in (\Sigma^*\times \Sigma^*)$ (dies ist gerade die Menge aller geordneten Paare von Wörtern über $\Sigma$) und bildet dieses auf ein neues Wort $xy\in\Sigma^*$ ab.
\end{Answer} 
\noindent
Die Verkettung $\text{Kon}$ über $\Sigma$ ist assoziativ über $\Sigma^*$, da offensichtlich
\begin{align*}
    \text{Kon}(x,\text{Kon}(y,z)) = x\cdot (y\cdot z) = xyz = (x\cdot y)\cdot z = \text{Kon}(\text{Kon}(x,y),z)
\end{align*}
gilt, für alle $x,y,z\in\Sigma^*$.\\

\noindent
Bei der Multiplikation in den reellen Zahlen ist die Zahl $1\in\mathbb{R}$ das neutrale Element, da $1\cdot x = x\cdot 1 = x$ für alle $x\in\mathbb{R}$ gilt.
\begin{Exercise}%[difficulty={1}]
\noindent
Welches ist das neutrale Element der Addition in den reellen Zahlen?\\
\vspace{0.6cm}
\end{Exercise}
\begin{Answer}
\noindent
Die Zahl $0\in\mathbb{R}$, da $0+x = x+0 = x$ für alle $x\in\mathbb{R}$.\\
\end{Answer}

\noindent
Für jedes $x \in\Sigma^*$ gilt
\begin{align*}
    x\cdot\lambda = \lambda\cdot x = x.
\end{align*}
Damit ist $\lambda$ das neutrale Element der Verkettung über $\Sigma^*$.\\

\begin{Exercise}%[difficulty={1}]
\noindent
Sei $x = a01b$ und $y = c21$. Bestimmen Sie $|x|$, $|y|$ und $|xy|$.\\
Seien allgemein $x,y\in\Sigma^*$ zwei Wörter für ein Alphabet $\Sigma$. Finden Sie einen Ausdruck für $|xy|$.\\
\vspace{0.6cm}
\end{Exercise}
\begin{Answer}
\noindent
Offensichtlich sind $|x| = |a01b|= 4$ und $|y| = |c21| = 3$ und somit $|xy| = |a01bc21| = 7$.\\

\noindent
Seien  $x,y\in\Sigma^*$ zwei Wörter für ein Alphabet $\Sigma$. Dann gilt\\
\begin{align*}
    |xy| = |x|+|y|
\end{align*}
\end{Answer} 

\begin{deff}[Umkehrung / Reversal]{deff:reversal}
Sei $n$ eine natürliche Zahl. Für ein Wort $x=x_1x_2...x_n$, mit $x_i\in\Sigma$ für $i\in\menge{1,2,...,n}$ bezeichnet $x^R = x_nx_{n-1}...x_1$ die \textbf{Umkehrung} (\textbf{Reversal}) von $x$. 
\end{deff}

\begin{Exercise}%[difficulty={1}]
\noindent
Sei $\Sigma$ ein Alphabet und $u,v\in\Sigma^*$ zwei Wörter. Beweisen oder widerlegen Sie die Aussage:
\begin{align*}
    (uv)^R = v^Ru^R
\end{align*}
\end{Exercise}
\begin{Answer}
\noindent
Da $u$ und $v$ Wörter sind, haben sie endliche Längen. Wir definieren $n$ und $m$ als $n =|u|$ und $m = |v|$. Dann hat $u$ die Form $u = u_1u_2...u_n$ und $v = v_1v_2...v_m$, wobei $u_1,u_2,...,u_n,v_1,v_2,...,v_m\in\Sigma$. Damit ist
\begin{align*}
    uv = u_1u_2...u_nv_1v_2...v_m
\end{align*}
und somit
\begin{align*}
    (uv)^R = v_m...v_2v_1u_n...u_2u_1 = v^Ru^R,
\end{align*}
da $v^R = v_m...v_2v_1$ und $u^R = u_n...u_2u_1$ gilt.
\clearpage
\end{Answer} 

\begin{deff}[Iteration eines Wortes]{deff:Iteration_Wort}
Sei $\Sigma$ ein Alphabet. Für alle $x\in\Sigma^*$ und alle $i\in\mathbb{N}$ definieren wir die $i$-te \textbf{Iteration} $x^i$ von $x$ als
\begin{align*}
    x^0 = \lambda,\quad x^1 = x, \quad x^i = xx^{i-1}.
\end{align*}
\end{deff}
\noindent
\underline{Beispiele}:\\
Sei $\Sigma = \menge{a,b,c}$. Wir können nun schreiben:
\begin{align*}
    aa &= a^2\\
    abababab &= (ab)^4\\
    cbbbbbab &= cb^5ab = c(bb)^2bab
\end{align*}

\noindent
Mit der folgenden Definition wollen wir den Begriff \textit{Teilwort}, für den wir eine gute Intuition haben, formalisieren. Ein Teilwort eines Wortes $x$ ist ein zusammenhängender Teil von $x$.
\begin{deff}[Teilwörter]{deff:Teilwort}
Seien $u,w\in\Sigma^*$ für ein Alphabet $\Sigma$.
\begin{itemize}
    \item $v$ heisst \textbf{Teilwort} von $w$ $\Leftrightarrow$ es existieren $x,y\in\Sigma^*$, so dass $w = xvy$
    \item $v$ heisst \textbf{Präfix} von $w$ $\Leftrightarrow$ es existiert $x\in\Sigma^*$, so dass $w = vx$
    \item $v$ heisst \textbf{Suffix} von $w$ $\Leftrightarrow$ es existiert $x\in\Sigma^*$, so dass $w = xv$
    \item $v\neq \lambda$ heisst \textbf{echtes} Teilwort (Präfix, Suffix) von $w$ $\Leftrightarrow$ $v\neq w$ und $v$ ein Teilwort (Präfix, Suffix) von $w$ ist
\end{itemize}
\end{deff}

\noindent
\underline{Beispiele}:\\
Sei $\Sigma = \menge{a,b,c}$. Das Wort $abc$ ist ein echtes Teilwort, echtes Präfix und ein echtes Suffix von $(abc)^3$. Das leere Wort $\lambda$ ist Teilwort von jedem Wort. Jedes Wort ist Teilwort von sich selbst (aber kein echtes Teilwort).

\begin{Exercise}%[difficulty={1}]
\noindent
(schwierig)\\
Sei $\Sigma = \menge{a_1,a_2,...,a_m}$ ein Alphabet der Grösse $m\in\mathbb{N}-\menge{0}$.
\begin{itemize}
    \item Welche Bedingungen muss ein Wort $x$ der Länge $n\leq m$ über $\Sigma$ erfüllen, damit es die maximale Anzahl von verschiedenen Teilwörtern hat?
    \item Wie viele verschiedene Teilwörter hat ein solches $x$?
\end{itemize}
\end{Exercise}
\begin{Answer}
\noindent
\begin{itemize}
    \item Die Bedingung ist, dass $x$ (Länge $n$) aus $n$ verschiedenen Buchstaben besteht (also aus lauter verschiedenen Buchstaben). Dann (und nur dann) sind offensichtlich alle möglichen Teilwörter paarweise verschieden und wir erreichen die maximale Anzahl von Teilwörtern.\\
    Ein solches Wort ist beispielsweise $w = a_1a_2...a_n$.
    \item Es genügt also, die Anzahl der Teilwörter von $w$ zu zählen. Wir werden die Anzahl der Teilwörter der Länge $i = 1, i = 2$ bis zur Länge $i = n$ zählen.\\
    \underline{$i = 1$}:
    \begin{align*}
        1:\quad &\blue{\underline{a_1}}a_2a_3a_4...a_{n-3}a_{n-2}a_{n-1}a_n\\
        2:\quad &a_1\blue{\underline{a_2}}a_3a_4...a_{n-3}a_{n-2}a_{n-1}a_n\\
        3:\quad &a_1a_2\blue{\underline{a_3}}a_4...a_{n-3}a_{n-2}a_{n-1}a_n\\
        &...\\
        (n-3):\quad &a_1a_2a_3a_4...\blue{\underline{a_{n-3}}}a_{n-2}a_{n-1}a_n\\
        (n-2):\quad &a_1a_2a_3a_4...a_{n-3}\blue{\underline{a_{n-2}}}a_{n-1}a_n\\
        (n-1):\quad &a_1a_2a_3a_4...a_{n-3}a_{n-2}\blue{\underline{a_{n-1}}}a_n\\
        n: \quad &a_1a_2a_3a_4...a_{n-3}a_{n-2}a_{n-1}\blue{\underline{a_n}}
    \end{align*}
    Offensichtlich gibt es $n$ viele Teilwörter der Länge $1$, nämlich genau die Teilwörter $\blue{a_1},\blue{a_2},...,\blue{a_n}$. Wir haben jeweils die Anfangsposition des Teilwortes in $w$ unterstrichen.
\end{itemize}
\underline{$i = 2$}:
    \begin{align*}
        1: \quad &\blue{\underline{a_1}a_2}a_3a_4...a_{n-3}a_{n-2}a_{n-1}a_n\\
        2: \quad &a_1\blue{\underline{a_2}a_3}a_4...a_{n-3}a_{n-2}a_{n-1}a_n\\
        3: \quad &a_1a_2\blue{\underline{a_3}a_4}...a_{n-3}a_{n-2}a_{n-1}a_n\\
        &...\\
        (n-3):\quad &a_1a_2a_3a_4...\blue{\underline{a_{n-3}}a_{n-2}}a_{n-1}a_n\\
        (n-2):\quad &a_1a_2a_3a_4...a_{n-3}\blue{\underline{a_{n-2}}a_{n-1}}a_n\\
        (n-1):\quad &a_1a_2a_3a_4...a_{n-3}a_{n-2}\blue{\underline{a_{n-1}}a_n}
    \end{align*}
    Es gibt $n-1$ viele Teilwörter der Länge $2$. Man beachte, dass sich die Anfangsposition der Teilwörter in $w$ um eine Position nach links verschoben hat, da das letzte Teilwort (ganz rechts) der Länge $2$ beim zweitletzten Symbol in $w$ beginnen muss.\\
\underline{$i = 3$}:
    \begin{align*}
        1: \quad &\blue{\underline{a_1}a_2}a_3a_4...a_{n-3}a_{n-2}a_{n-1}a_n\\
        2: \quad &a_1\blue{\underline{a_2}a_3a_4}...a_{n-3}a_{n-2}a_{n-1}a_n\\
        &...\\
        (n-3):\quad &a_1a_2a_3a_4...\blue{\underline{a_{n-3}}a_{n-2}a_{n-1}}a_n\\
        (n-2):\quad &a_1a_2a_3a_4...a_{n-3}\blue{\underline{a_{n-2}}a_{n-1}a_n}
    \end{align*}
    Es gibt $n-2$ viele Teilwörter der Länge $3$.\\
    
    \noindent
    Es ist nun nicht mehr schwierig zu sehen, dass es für ein Teilwort der Länge $i\in\menge{1,2,...,n}$ genau $n-i+1$ mögliche Anfangspositionen in $w$ gibt, da die Positionen $n-i+2,3,...,n$ nicht möglich sind (\quotes{zu wenig Platz}). Nun summieren wir die Anzahl aller unterschiedlichen (nicht leeren) Teilwörter von $w$ der Längen $i=1,2,...n$ auf:
    \begin{align*}
        \sum_{i=1}^n (n-i+1) = n(n+1)-\sum_{i=1}^n i = n(n+1) - \frac{n(n+1)}{2}=\frac{n(n+1)}{2}
    \end{align*}
    Hinzu kommt noch das leere Wort (welches Teilwort jedes Wortes ist) und als Resultat finden wir, dass es
    \begin{align*}
        \frac{n(n+1)}{2} + 1
    \end{align*}
    viele verschiedene Teilwörter von $w$ gibt.
\end{Answer}
\clearpage
\noindent
Wir wollen an dieser Stelle noch eine weitere nützliche Schreibweise definieren.\\

\noindent
Sei $x\in\Sigma^*$ und $a\in\Sigma$, dann ist $|x|_a$ definiert als die Anzahl der Vorkommen von $a$ in $x$.\\
\underline{Beispiel}:\\
Sei $v = babaac$, dann ist $|v|_a = 3$, $|v|_b = 2$ und $|v|_c = 1$. Offensichtlich gilt für alle $x\in\Sigma^*$
\begin{align*}
    |x| = \sum_{a\in\Sigma}|x|_a
\end{align*}

\noindent
Nun kommen wir zur Definition einer Sprache. Dies wird für uns einer der wichtigsten Begriffe sein.
\begin{deff}[Sprache]{deff:Sprache}
Eine \textbf{Sprache} $L$ über einem Alphabet $\Sigma$ ist eine Teilmenge von $\Sigma^*$. Das Komplement $L^C$ der Sprache $L$ bezüglich $\Sigma$ ist die Sprache $\Sigma^*-L$.\\

\noindent
$L_{\emptyset} = \emptyset$ ist die leere Sprache\\
$L_{\lambda} = \menge{\lambda}$ ist die einelementige Sprache, die nur aus dem leeren Wort besteht.\\

\noindent
sind $L_1$ und $L_2$ Sprachen über $\Sigma$, so bezeichnet
\begin{align*}
    L_1\cdot L_2 = L_1L_2 = \mengecm{vw}{v\in L_1\; \text{und}\; w\in L_2}
\end{align*}
die \textbf{Konkatenation} von $L_1$ und $L_2$.\\

\noindent
Ist $L$ eine Sprache über $\Sigma$, so definieren wir die Iterationen
\begin{align*}
    L^0 &= L_{\lambda}, L^{i+1} = L^i\cdot L\;\text{für alle}\; i\in\mathbb{N},\\
    L^* &= \bigcup_{i\in\mathbb{N}}L^i\;\text{und}\; L^+ = \bigcup_{i\in\mathbb{N}}L^i\cdot L
\end{align*}
\end{deff}
$L^*$ nennt man den \textbf{Kleene'schen Stern}\footnote{benannt nach dem US-Amerikanischen Mathematiker Stephen Cole Kleene} von $L$.\\


\clearpage
\noindent
\underline{Beispiele}:\\
Die folgenden Mengen sind Beispiele von Sprachen über $\Sigma = \menge{a,b}$.
\begin{itemize}
    \item $L_1 = \emptyset$
    \item $L_2 = \menge{\lambda}$
    \item $L_3 = \Sigma^* = \menge{\lambda, a,b,aa,...}$
    \item $L_4 = \Sigma^+ = \menge{a,b,aa,...}$
    \item $L_5 = \Sigma$
    \item $L_6 = \mengecm{a^p}{a\;\text{ist eine Primzahl}}$
    \item $L_5 = \menge{a}^* = \menge{\lambda, a, aa, aaa, aaaa, ...}$
    \item $\Sigma^2 = \menge{aa,ab,ba,bb}$
\end{itemize}
Man beachte, dass $\Sigma^i = \mengecm{x\in\Sigma^*}{|x| = i}$ für eine natürliche Zahl $i$, und dass $L_{\emptyset}\cdot L = \emptyset$, $L_{\lambda}\cdot L = L$.\\

\noindent
Die Menge aller grammatisch korrekten englischen Texte ist eine Sprache über $\Sigma_{\text{Tastatur}}$.\\
Die Menge aller syntaktisch korrekten Programme in {\tt C++} ist auch eine Sprache über $\Sigma_{\text{Tastatur}}$.

\begin{Exercise}%[difficulty={1}]
\noindent
Sei $L_1 = \menge{\lambda, a, b}$ und sei $L_2 = \menge{a^4, a^2b}$. Welche Wörter Liegen in der Sprache $L_3 = L_1L_2$?
\end{Exercise}
\begin{Answer}
\noindent
\begin{align*}
    L_3 = \menge{a^4, a^2b, a^5, a^3b,ba^4, ba^2b}
\end{align*}
\end{Answer}

\begin{Exercise}%[difficulty={1}]
\noindent
Sei $k\in\mathbb{N}$. Geben Sie ein Alphabet $\Sigma$ und zwei Sprachen $L_1$ und $L_2$ über $\Sigma$ an, sodass
\begin{align*}
    |L_1| = k\quad \text{und} \quad |L_1L_2| = k + 1.
\end{align*}
\end{Exercise}
\begin{Answer}
\noindent
Wir wählen $\Sigma = \menge{0}$ und
\begin{align*}
    L_1 = \menge{0,00,...,0^k} = \mengecm{0^i}{1\leq i \leq k}
\end{align*}
und
\begin{align*}
    L_2 = \menge{\lambda,0}.
\end{align*}
Dann gilt
\begin{align*}
    L_1\cdot L_2 &= L_1\cdot \menge{\lambda}\cup L_1\cdot \menge{0}\\
    &= L_1 \cup \mengecm{0^i0}{1\leq i\leq k}\\
    &= L_1 \cup \mengecm{0^i}{2\leq i \leq k+1}\\
    &= \mengecm{0^i}{1\leq i\leq k+1},
\end{align*}
und somit $|L_1L_2| = k+1$.
\end{Answer}